% Part: first-order-logic
% Chapter: natural-deduction
% Section: quantifier-rules

\documentclass[../../../include/open-logic-section]{subfiles}

\begin{document}

\olfileid{fol}{ntd}{qrl}

\olsection{قواعد المكممات}

\subsection{قواعد~$\lforall$}

\begin{defish}
\AxiomC{$!A(a)$}
\RightLabel{\Intro{\lforall}}
\UnaryInfC{$\lforall[x][\Atom{!A}{x}]$}
\DisplayProof
\hfill
\AxiomC{$\lforall[x][\Atom{!A}{x}]$}
\RightLabel{\Elim{\lforall}}
\UnaryInfC{$!A(t)$}
\DisplayProof
\end{defish}

في قاعدتي~$\lforall$، يكون~$t$ حدًا مغلقًا (أي حدًا لا يحتوي على أي
متغيرات)، ويكون~$a$ !!a{constant} لا يرد في النتيجة
$\lforall[x][!A(x)]$، ولا في أي افتراض !!{undischarged} في
ال!!{derivation} المنتهي بالمقدمة~$!A(a)$. ونسمي~$a$
\emph{المتغير المميَّز} لاستدلال~\Intro{\lforall}.
\footnote{نستعمل مصطلح ``المتغير المميَّز'' حتى وإن كان~$a$ في القاعدة
أعلاه رمزًا ثابتًا، وذلك لأسباب تاريخية.}

\subsection{قواعد~$\lexists$}

\begin{defish}
\AxiomC{$\Atom{!A}{t}$}
\RightLabel{\Intro{\lexists}}
\UnaryInfC{$\lexists[x][\Atom{!A}{x}]$}
\DisplayProof
\hfill
\AxiomC{$\lexists[x][\Atom{!A}{x}]$}
\AxiomC{[$\Atom{!A}{a}$]$^n$}
\DeduceC{$!C$}
\DischargeRule{\Elim{\lexists}}{n}
\BinaryInfC{$!C$}
\DisplayProof
\end{defish}

مرة أخرى، يكون~$t$ حدًا مغلقًا، ويكون~$a$ !!a{constant} لا يرد في
المقدمة~$\lexists[x][!A(x)]$، ولا في النتيجة~$!C$، ولا في أي افتراض
!!{undischarged} في ال!!{derivation}s المنتهية بالمقدمتين (عدا
الافتراضات~$!A(a)$). ونسمي~$a$ \emph{المتغير المميَّز} لاستدلال
\Elim{\lexists}.

يسمى الشرط القاضي بألا يرد المتغير المميَّز في النتائج ولا في أي
افتراض !!{undischarged} في ال!!{derivation}s المؤدية إلى مقدمتي استدلال
\Intro{\lforall} أو~\Elim{\lexists} \emph{شرط المتغير المميَّز}.

\begin{explain}
تذكّر الاصطلاح القائل إننا عندما تكون~$!A$ !!a{formula} يكون فيها
ال!!{variable}~$x$ حرًا، نشير إلى ذلك بكتابة~$!A(x)$. وفي السياق نفسه،
تكون~$!A(t)$ اختصارًا لـ~$\Subst{!A}{t}{x}$. ولذلك يمكننا أيضًا كتابة
قاعدة~$\Intro\lexists$ هكذا:
\begin{prooftree}
\AxiomC{$\Subst{!A}{t}{x}$}
\RightLabel{\Intro{\lexists}}
\UnaryInfC{$\lexists[x][!A]$}
\end{prooftree}
لاحظ أن~$t$ قد يرد أصلًا في~$!A$؛ فمثلًا قد تكون~$!A$ هي
$\Atom{\Obj P}{t,x}$. ومن ثم فإن استنتاج
$\lexists[x][\Atom{\Obj P}{t,x}]$ من~$\Atom{\Obj P}{t,t}$ تطبيق
صحيح لـ~$\Intro\lexists$---إذ يجوز لك أن ``تستبدل'' بموضع واحد أو
أكثر، وليس بالضرورة كل المواضع، للحد~$t$ في المقدمة ال!!{variable}
المقيد~$x$. غير أن شرطي المتغير المميَّز في~$\Intro\lforall$
و$\Elim\lexists$ يقتضيان ألا يرد ال!!{constant}~$a$ في~$!A$.
ولذلك لا يمكنك أن تستنتج على نحو صحيح
$\lforall[x][\Atom{\Obj P}{a,x}]$ من~$\Atom{\Obj P}{a,a}$ باستعمال
$\Intro\lforall$.
\end{explain}

\begin{explain}
لا توجد قيود في~\Intro{\lexists} و\Elim{\lforall}، ويمكن أن يكون
الحد~$t$ أي شيء، فلا داعي إلى القلق من أي شروط. أما في قاعدتي
\Elim{\lexists} و\Intro{\lforall}، فيقتضي شرط المتغير المميَّز ألا
يرد ال!!{constant}~$a$ في أي موضع من النتيجة أو من افتراض
!!{undischarged}. والشرط ضروري لضمان سلامة النسق، أي إنه لا
!!{derive}s إلا !!{sentence}s تلزم عن الافتراضات ال!!{undischarged}
التي تشتق منها. ومن دون هذا الشرط، يكون الآتي جائزًا:
\begin{prooftree}
  \AxiomC{$\lexists[x][!A(x)]$}
  \AxiomC{$\Discharge{!A(a)}{1}$}
  \RightLabel{*\Intro{\lforall}}
  \UnaryInfC{$\lforall[x][!A(x)]$}
  \RightLabel{\Elim{\lexists}}
  \BinaryInfC{$\lforall[x][!A(x)]$}
\end{prooftree}
لكن~$\lexists[x][!A(x)] \Entails/ \lforall[x][!A(x)]$.

بما أن قواعد الحذف للمكممات لا تسمح إلا بإحلال حدود مغلقة محل
ال!!{variable}s، فإنه يلزم أن تكون أي !!{formula} يمكن اشتقاقها من
مجموعة من ال!!{sentence}s هي نفسها !!a{sentence}.
\end{explain}


\end{document}
